\section{Discussion and Conclusion}
\label{sec:discussion}

\subsection{What a \knowact Result Would Mean}

A benchmark score is evidence only relative to the explanations excluded by the design.

Low reconstruction error may reflect dialogue inference, a population prior, leaked information, or a
good guess.

\knowact strengthens interpretation through controlled access, evidence references, and non-adaptive
comparisons.

The resulting claims form a hierarchy:

\begin{enumerate}[leftmargin=1.7em]
    \item \emph{reconstruction}: the final map matches the reviewed map;
    \item \emph{active acquisition}: adaptive policies outperform matched non-adaptive controls;
    \item \emph{state use}: working-map ablations identify a mechanism for that advantage;
    \item \emph{graph-aware diagnosis}: graph structure improves efficiency without unsupported propagation;
    \item \emph{proxy robustness}: the advantage persists across simulator configurations;
    \item \emph{human transfer}: the advantage persists in matched human interaction.
\end{enumerate}

Evidence for a lower level does not entail a higher one.

Accurate reconstruction is not active user modeling if a fixed script achieves the same result. Failure
to transfer to humans indicts the simulator proxy, not necessarily task consistency.

The interaction is active information acquisition over a structured latent state.

It differs from pool-based active learning because probes are open utterances that require grounding. It
differs from adaptive testing because the item bank is neither fixed nor psychometrically calibrated.

This expressivity makes provenance, grounding, and proxy validation part of the measurement argument.

Even a fully validated result would remain diagnostic rather than task-complete.

\knowact asks whether an agent acquires a useful user model. It does not establish improved teaching,
recommendation, or collaboration.

Connecting reconstruction to downstream benefit is a separate claim and a natural next benchmark layer.

\subsection{Limitations and Threats to Validity}

\paragraph{The target is authored, not discovered.}
The graph and rubrics determine which distinctions count as knowledge.

Expert review and source provenance make that choice inspectable, not uniquely correct. Rankings may
change under another granularity, relation set, or rubric.

\paragraph{Simulator fidelity.}
A simulator may be consistent with a reviewed map while expressing partial knowledge unlike a person.

Context isolation limits raw-field exposure, but does not establish semantic safety or behavioral
realism. Held-out answers and matched human studies must test state fidelity and ranking transfer.

\paragraph{Artifact dependence.}
The graph, rubrics, maps, and simulator jointly operationalize the construct. Shared design choices may
favor some model families.

Source grounding, deterministic fixtures, review, simulator variation, and human validation reduce but
cannot remove this dependence.

\paragraph{Simplified knowledge state.}
A static L0--L5 state cannot capture context, procedural fluency, confidence, forgetting, or learning.

Freezing state makes reconstruction identifiable, but it treats explanations as observations rather
than teaching events. Dynamic diagnosis needs a transition model and different causal interpretation.

\paragraph{Initial domain scope.}
The first release focuses on statistical learning and a small authored graph. It can validate the
protocol, not cross-domain generality.

Additional domains must vary source style, topology, concept granularity, and expressible misconceptions.

\paragraph{Operational claims.}
Success would show behavior compatible with structured user modeling under benchmark constraints.

It would not establish human-like beliefs, consciousness, or a faithful psychological mechanism. The
term functional ToM names the tested role, not an ontological status.

\subsection{Conclusion and Outlook}

Evaluating user-aware agents requires more than placing them in a conversation.

The evaluator must know the user state at stake, make evidence acquisition consequential, and separate
accurate inference from leakage, priors, and generic interaction skill.

\knowact turns a reviewed knowledge map into an explicit latent target and lets agents choose probes
under a budget.

It exposes the link from visible evidence to working state and action. The \sagesim boundary and
deterministic score make the loop auditable without equating simulation with human interaction.

The immediate next step is to test, not broaden, the claim.

Adaptive controls, mechanism ablations, simulator variation, and human episodes must establish which
evidence level \knowact supports.

Only then should the benchmark extend to dynamic learning, downstream assistance, and new domains.
